% !TEX root = ../../hw05.tex
\begin{proof}[Solution]
  Let $\varepsilon>0$. Our goal is to make
  $|g(f(x))-g(f(x_0))|<\varepsilon$ by taking $x$ sufficiently close
  to $x_0$. We first find how close $f(x)$ must be to $f(x_0)$.

  Write $y_0=f(x_0)$. Since $g$ is continuous at $y_0$, there is an
  $\eta>0$ such that
  \[
    |y-y_0|<\eta
      \quad\Longrightarrow\quad |g(y)-g(y_0)|<\varepsilon.
  \]
  Since $f$ is continuous at $x_0$, there is a $\delta>0$ such that
  \[
    |x-x_0|<\delta
      \quad\Longrightarrow\quad |f(x)-f(x_0)|<\eta.
  \]
  The choices are made in this order because the tolerance $\eta$ for
  $f$ comes from the continuity of $g$.
  \autoref{fig:hw05:composition-neighbourhoods} illustrates the two steps.

  \begin{marginfigure}
    \centering
    % !TEX root = ../../problems/hw05.tex
% These neighbourhoods are schematic and do not have a common scale.
% Shared physical dimensions for interval and number-line diagrams.
\tikzset{
  interval diagram/.style={
    x=1cm,y=6mm,font=\footnotesize,
    color=black,line width=0.5pt,line cap=round,line join=round,
    >={Stealth[length=4pt,width=3pt]}
  },
  interval endpoint/.style={
    circle,draw=black,line width=0.5pt,
    minimum size=4pt,inner sep=0pt,outer sep=0pt
  },
  interval open/.style={interval endpoint,fill=white},
  interval closed/.style={interval endpoint,fill=black},
  interval label/.style={inner sep=1pt}
}

\begin{tikzpicture}[interval diagram,x=10mm,y=1mm]
  \foreach \height/\centre/\tolerance in {
    0/{x_0}/{\delta},-24/{y_0=f(x_0)}/{\eta},
    -48/{g(y_0)}/{\varepsilon}}{
    \fill[gray!10] (-1.25,{\height-2}) rectangle (1.25,{\height+2});
    \draw[<->,gray] (-1.6,\height) -- (1.6,\height);
    \node[interval open] at (-1.25,\height) {};
    \node[interval open] at (1.25,\height) {};
    \fill (0,\height) circle[radius=1.6pt];
    \node[below=5pt] at (0,\height) {$\centre$};
    \draw[<->] (0,{\height+7}) -- (1.25,{\height+7})
      node[midway,above=4pt] {$\tolerance$};
  }
  % Leave a separate gutter for arrows, away from the centre labels.
  \draw[->] (-1.85,-4) -- (-1.85,-20)
    node[midway,left=4pt] {$f$};
  \draw[->] (-1.85,-28) -- (-1.85,-44)
    node[midway,left=4pt] {$g$};
\end{tikzpicture}

    \caption[Neighbourhoods for a composition.]{%
      Schematic neighbourhoods for a composition.
      First choose $\eta$ to reach the required output tolerance $\varepsilon$.
      Then choose $\delta$ to reach that intermediate tolerance $\eta$.
      The arrows show the order in which inputs pass through the maps.}
    \label{fig:hw05:composition-neighbourhoods}
  \end{marginfigure}

  If $|x-x_0|<\delta$, the second implication places $y=f(x)$ within
  $\eta$ of $y_0$. Applying the first implication with this value of $y$
  gives
  \[
    |h(x)-h(x_0)|
      =|g(f(x))-g(f(x_0))|<\varepsilon.
  \]
  This is the definition of continuity of $h=g\circ f$ at $x_0$.
\end{proof}
